\section{特殊函数}

在本附录中，我们将在不同边界条件下，研究受限连续介质中单粒子态的渐近分布。为简化起见，我们考虑一个边长分别为 \(a\)、\(b\) 和 \(c\) 的长方体封闭区域。自由粒子薛定谔方程的可接受解为
\begin{equation}\tag{17.1.1}\label{eq:17.1.1}
\nabla ^{2} \psi + k ^{2} \psi = 0  \left( k = p \hbar ^{- 1} \right) ,
\end{equation}
这些解满足狄利克雷边界条件（即，在边界上处处 \(\psi=0\)），其形式为
\begin{equation}\tag{17.1.2}\label{eq:17.1.2}
\psi _ { l m n } ( \pmb { r } ) \propto \sin \left( \frac { l \pi x } { a } \right) \sin \left( \frac { m \pi y } { b } \right) \sin \left( \frac { n \pi z } { c } \right) ,
\end{equation}
其中
\begin{equation}\tag{17.1.3}\label{eq:17.1.3}
k = \pi \left( \frac { l ^{2} } { a ^{2} } + \frac { m ^{2} } { b ^{2} } + \frac { n ^{2} } { c ^{2} } \right) ^{1 / 2} ;  l , m , n = 1 , 2 , 3 , \ldots .
\end{equation}
需要注意的是，在这种情况下，量子数 \(l\)、\(m\) 或 \(n\) 都不能为零，因为那样会导致波函数恒等于零。相反，若我们采用诺伊曼边界条件（即，在边界上处处 \(\frac{\partial\psi}{\partial n}=0\)），则得到的解为
\begin{equation}\tag{17.1.4}\label{eq:17.1.4}
\psi _ { l m n } ( \pmb { r } ) \propto \cos \left( \frac { l \pi x } { a } \right) \cos \left( \frac { m \pi y } { b } \right) \cos \left( \frac { n \pi z } { c } \right) ,
\end{equation}
其中
\begin{equation}\tag{17.1.5}\label{eq:17.1.5}
l , m , n = 0 , 1 , 2 , \ldots ;
\end{equation}
显然，此时量子数的零值已获准！然而，在每一种情况下，量子数的负整数值都不会产生新的波函数。
总共有 \(g(K)\) 个不同的波函数 \(\psi\)，其波数 \(k\) 不超过给定的值 \(K\)，可以表示为
\begin{equation}\tag{17.1.6}\label{eq:17.1.6}
g ( K ) = \sum _ { l , m , n } ^{\prime} f ( l , m , n ) ,
\end{equation}
其中对于属于集合（3）或（5）的 \((l,m,n)\) 数字，\(f(l,m,n)=1\)；每种情况下的求和 \(\sum'\) 均受以下条件限制：
\begin{equation}\tag{17.1.7}\label{eq:17.1.7}
\left( \frac { l ^{2} } { a ^{2} } + \frac { m ^{2} } { b ^{2} } + \frac { n ^{2} } { c ^{2} } \right) \leq \frac { K ^{2} } { \pi ^{2} } .
\end{equation}
我们现在定义一个求和
\begin{equation}\tag{17.1.8}\label{eq:17.1.8}
G ( K ) = \sum _ { l , m , n } ^{\prime} f ^{*} ( l , m , n ) ,
\end{equation}
其中对于所有整数值 \(l\)、\(m\) 和 \(n\)（正整数、负整数或零），\(f^{*}(l,m,n)=1\)，而 \((l,m,n)\) 数字的限制条件与（7）中所述一致。通过建立项之间的对应关系，我们可以证明：
\begin{equation}\tag{17.1.9}\label{eq:17.1.9}
\begin{aligned}
\sum_{l,m,n}^{\prime} f(l,m,n)
&= \frac{1}{8}\Bigg[
\sum_{l,m,n}^{\prime} f^{*}(l,m,n) \\
&\mp \Big\{
\sum_{l,m}^{\prime} f^{*}(l,m,0)
+ \sum_{l,n}^{\prime} f^{*}(l,0,n)
+ \sum_{m,n}^{\prime} f^{*}(0,m,n)
\Big\} \\
&+ \Big\{
\sum_{l}^{\prime} f^{*}(l,0,0)
+ \sum_{m}^{\prime} f^{*}(0,m,0)
+ \sum_{n}^{\prime} f^{*}(0,0,n)
\Big\}
\mp 1
\Bigg].
\end{aligned}
\end{equation}
这里的正号与负号分别对应狄利克雷边界条件与诺伊曼边界条件。
显然，（9）右侧的第一个求和表示椭球体中的晶格点数量（\(^{1}\)）——即 \(X^{2}/a^{2}+Y^{2}/b^{2}+Z^{2}/c^{2}\) = \(K^{2}/\pi^{2}\)；接下来的三个求和分别表示该椭球体在以 \(Z\)、\(Y\) 和 \(X\) 平面为截面的椭圆中的晶格点数量；最后三个求和则表示椭球体主轴上的晶格点数量。如果 \(a\)、\(b\) 和 \(c\) 相较于 \(\pi/K\) 而言足够大，我们就可以将这些数字分别替换为相应的体积、面积和长度，其结果为
\begin{equation}\tag{17.1.10}\label{eq:17.1.10}
\begin{aligned}
g(K)={}&\frac{K^3}{6\pi^2}(abc)
\mp \frac{K^2}{8\pi}(ab+ca+bc) \\
&+ \frac{K}{4\pi}(a+b+c)
\mp \frac{1}{8}
+ E(K).
\end{aligned}
\end{equation}
此处的"\(E(K)\)"一词，指在进行上述替换操作时所造成的净误差。因此，我们发现，本结果中的主要项与封闭区域的体积成正比；而第一修正项则与该区域的表面积成正比（因而代表一种"表面效应"）；接下来的阶数项则以"边缘效应"和"角落效应"的形式出现。
如今，查阅有关"给定区域内晶格点数量"计算方法的相关文献可知，式（10）中误差项\(E(K)\)的阶为\(O(K^{\alpha})\)，其中\(1 < \alpha < 1.4\)；因此，\(g(K)\)的式（10）仅在表面项这一层次上才具有可靠性。鉴于此，我们可以写成：
\begin{equation}\tag{17.1.11}\label{eq:17.1.11}
g ( K ) = \frac { K ^{3} } { 6 \pi ^{2} } V \mp \frac { K ^{2} } { 16 \pi } S + \mathrm { a ~ l o w e r - o r d e r ~ r e m a i n d e r } ;
\end{equation}
¹ 术语"在椭球体内"指的是"不位于椭球体外部"，也就是说，那些"位于椭球体上的"晶格点同样被纳入考量。后续各处的类似表述亦具有相似的含义。
以\(\varepsilon^{*}\)为变量，其中
\begin{equation}\tag{17.1.12}\label{eq:17.1.12}
\varepsilon ^{*} = \frac { 8 m L ^{2} } { h ^{2} } \varepsilon = \frac { 4 L ^{2} } { h ^{2} } P ^{2} = \frac { L ^{2} } { \pi ^{2} } K ^{2} ,
\end{equation}
式（11）可简化为文中的式（1.4.15）和式（1.4.16）。
在周期性边界条件的情况下，即
\begin{equation}\tag{17.1.13}\label{eq:17.1.13}
\psi ( x , y , z ) = \psi ( x + a , y , z ) = \psi ( x , y + b , z ) = \psi ( x , y , z + c ) ,
\end{equation}
合适的波函数为
\begin{equation}\tag{17.1.14}\label{eq:17.1.14}
\psi _ { l m n } ( \mathbf { r } ) \propto \exp \{ i ( \mathbf { k } \cdot \mathbf { r } ) \} ,
\end{equation}
其中
\begin{equation}\tag{17.1.15}\label{eq:17.1.15}
k = 2 \pi \left( { \frac { l } { a } } , { \frac { m } { b } } , { \frac { n } { c } } \right) ;  l , m , n = 0 , \pm 1 , \pm 2 , \ldots .
\end{equation}
自由粒子态的数量\(g(K)\)现可表示为
\begin{equation}\tag{17.1.16}\label{eq:17.1.16}
g ( K ) = \sum _ { l , m , n } ^{\prime} f ^{*} ( l , m , n ) ,
\end{equation}
使得
\begin{equation}\tag{17.1.17}\label{eq:17.1.17}
( l ^{2} / a ^{2} + m ^{2} / b ^{2} + n ^{2} / c ^{2} ) \leq K ^{2} / ( 4 \pi ^{2} ) .
\end{equation}
这正是半轴长度分别为\(Ka/2\pi\)、\(Kb/2\pi\)和\(Kc/2\pi\)的椭球体内的晶格点数量；若考虑先前情形中所做的近似，其值恰好等于式（11）中的体积项。因此，在周期性边界条件的情况下，我们在状态密度的表达式中并未得到表面项。
如需进一步了解该主题的相关信息，请参阅费多索夫（1963年、1964年）、帕特里亚（1966年）、查巴与帕特里亚（1973年）以及巴尔特斯与希尔夫（1976年）的著作。

